% Chapter 4: Higher-order and one-sided stencils (methods author). PLAN.md section 1.4; sidebars S6 and S7;
% figures F6 to F9; examples E5 and E6; the margin table of node counts (style decision 1).
% Revised at stage 6 per reviews/R1-decision.md (T6, T8, T10, numerics-08, numerics-17b, numerics-17d, evidence-14,
% evidence-18, evidence-22, library-13, engineer-19, engineer-20); the changes are listed in responses/R1-04-stencils.md.
\chapter{Higher-order and one-sided stencils}\label{ch:stencils}

\Cref{ch:taylor} built the forward and the central difference from Taylor's theorem, and \cref{ch:rounding} showed
that the central difference reaches about \num{\valHeadlineFirstDigits} digits in double and no more, because below
its best step the rounding of $f$ takes over. This chapter asks what more points buy. The answer is a higher accuracy
order, which lets the step grow and the error fall, at a price in rounding, in smoothness and in evaluations. It
also answers a practical question that the central difference cannot: how to differentiate at the edge of a domain,
where $f$ is defined on one side of $x$ only. The thread continues: $\cp = \dd \enth / \dd T$ of nitrogen at
$T_0 = \SI{1000}{\kelvin}$, and at the lower end of the fit's range.

\section{Stencils: nodes, weights and the moment conditions}\label{sec:stencils-moments}

Every rule of this book has the same shape. It samples $f$ at a few points $x_k = x + n_k h$ and combines the
samples with fixed weights:
\begin{equation}
  f^{(\dorder)}(x) \approx \frac{1}{h^{\dorder}} \sum_k w_k\, f(x + n_k h).
  \label{eq:stencils-rule}
\end{equation}
\mnote{A \emph{node} is an integer offset $n_k$; the \emph{sample} is $f_k = f(x + n_k h)$; the \emph{stencil} is
the pattern of nodes; a \emph{rule} is a stencil with its weights.}The integers $n_k$ are the nodes, the rational
numbers $w_k$ the weights, and $\dorder$ the derivative order, which is 1 in this chapter. The central difference
has the nodes $-1, 1$ with the weights $-\tfrac12, \tfrac12$; the forward difference the nodes $0, 1$ with the
weights $-1, 1$. Where do the weights come from, and how are more of them found?

Taylor's theorem answers both. Expanding each sample about $x$,
$f(x + n_k h) = \sum_j f^{(j)}(x)\, n_k^j h^j / j!$, and collecting the powers of $h$ gives
\begin{equation}
  \sum_k w_k\, f(x + n_k h) = \sum_{j \ge 0} \frac{\mom{j}}{j!}\, h^j f^{(j)}(x),
  \qquad \mom{j} = \sum_k w_k\, n_k^j .
  \label{eq:stencils-moments}
\end{equation}
The numbers $\mom{j}$ are the moments of the rule, and they decide everything. The weighted sum is $h^{\dorder}
f^{(\dorder)}$ plus an error of order $h^{\dorder + \aorder}$ exactly when $\mom{\dorder} = \dorder!$ and every
other moment up to $\mom{\dorder + \aorder - 1}$ vanishes; these are the moment conditions. The first moment that
does not vanish, $\mom{\dorder + \aorder}$, gives the leading error term of the rule,
$\errc\, h^{\aorder} f^{(\dorder + \aorder)}$ with $\errc = \abs{\mom{\dorder + \aorder}} / (\dorder + \aorder)!$,
and $\aorder$ is the accuracy order.

Take the four-point central rule for $f'$, the library's \code{Central<4>}. Its nodes are $-2, -1, 1, 2$. A rule
that is to be exact for the odd function $f(x + t) = t$ and give zero for even functions has antisymmetric
weights, $w_{-1} = -a$, $w_1 = a$, $w_{-2} = -b$, $w_2 = b$, so every even moment vanishes by itself. The two
odd moments that remain give two conditions, $\mom{1} = 2a + 4b = 1$ and $\mom{3} = 2a + 16b = 0$, whose solution
is $a = \tfrac23$ and $b = -\tfrac{1}{12}$:
\begin{equation}
  f'(x) = \frac{f(x - 2h) - 8 f(x - h) + 8 f(x + h) - f(x + 2h)}{12 h} + \frac{h^4}{30} f^{(5)}(\xi).
  \label{eq:stencils-central-four}
\end{equation}
The error comes from $\mom{5} = 2a + 64 b = -4$, so $\errc = 4/5! = 1/30$, and the rule's value is $f'$ minus
that term: it underestimates $f'$ where $f^{(5)} > 0$.\sidenote{The single intermediate point $\xi$ in the span
of the stencil comes from the interpolation error formula (\cref{sec:derivations-moments}); the next term is $h^6$,
since all even moments vanish.} Two more points have raised the order from 2 to 4. \Cref{tab:weights-central-first}
in \cref{ch:tables} lists the nodes, weights, $\ampl$ and $\errc$ of every central rule for $f'$, and the general
machinery that produces them, one formula over the nodes, is in \cref{sb:moments}. One consequence of the moment
conditions is worth keeping in mind: a rule of order $\aorder$ is exact for every polynomial of degree up to
$\dorder + \aorder - 1$, because all the derivatives in its error terms vanish.

\begin{margintable}
\caption{Nodes and evaluations of \code{fn} per call of a stencil, for the derivative order $\dorder$. The central
rules omit node 0 for odd $\dorder$. Ridders' nodes per level are in \cref{tab:api-methods}.\label{tab:stencils-nodes}}
\centering\scriptsize
\setlength{\tabcolsep}{2.5pt}
\begin{tabular}{@{}lcc@{}}
\toprule
Method & Nodes & Evaluations\\
\midrule
\code{Central<P>} & $\aorder + \dorder - 1$ & $\aorder + \dorder - 1$\\
\code{Forward<P>} & $\aorder + \dorder$ & $\aorder + \dorder$\\
\code{Backward<P>} & $\aorder + \dorder$ & $\aorder + \dorder$\\
\bottomrule
\end{tabular}
\end{margintable}

\begin{sidebar}[label=sb:moments]{The moment conditions and the Lagrange generator}
A rule for the $\dorder$-th derivative of accuracy order $\aorder$ on the nodes $n_1, \ldots, n_m$ satisfies
\[
  \mom{j} = \sum_k w_k n_k^j = 0 \quad \text{for } j \ne \dorder,\ 0 \le j \le \dorder + \aorder - 1,
  \qquad \mom{\dorder} = \dorder!,
\]
a linear system for the weights. The library does not solve it as a system. It uses the Lagrange formula
(Burden et al.\cite[Sect.~4.1]{burden2016}): the polynomial $p$ that interpolates the samples at the nodes, in the
scaled variable $t$ with $x + t h$, is $p(t) = \sum_k f_k L_k(t)$ with $L_k(t) = \prod_{j \ne k} (t - n_j) / (n_k - n_j)$,
and its $\dorder$-th derivative at $t = 0$ gives the weights directly:
\[
  w_k = L_k^{(\dorder)}(0) = \dorder!\,[t^{\dorder}] \prod_{j \ne k} \frac{t - n_j}{n_k - n_j},
\]
where $[t^{\dorder}]$ takes the coefficient of $t^{\dorder}$. Why this solves the moment conditions: $p$
reproduces every polynomial of degree below the number of nodes, so for $f = t^j$ the weighted sum equals the
$\dorder$-th derivative of $t^j$ at 0, which is $\dorder!$ for $j = \dorder$ and 0 otherwise
(\cref{sec:derivations-lagrange}). The library's own worked example is the node $-1$ of \code{Central<4>}:
$\prod_{j \ne k} (t - n_j) = (t + 2)(t - 1)(t - 2) = t^3 - t^2 - 4t + 4$, whose $t^1$ coefficient is $-4$, over
$\prod (n_k - n_j) = (-1 + 2)(-1 - 1)(-1 - 2) = 6$, so $w = -2/3$, as in \cref{eq:stencils-central-four}.

The library stores every rule as integer numerators $c_k$ over one common denominator $\den$, so that
$w_k = c_k / \den$; for \code{Central<4>} the numerators are $1, -8, 8, -1$ over 12, and with the integers the
moment conditions read $\mom{\dorder} = \dorder!\,\den$. This is the form \cref{ch:tables} tabulates, and it keeps
the weights exact in every real type. One \code{consteval} function evaluates the Lagrange formula for every
family and order; each family only supplies its nodes. It works in exact rational arithmetic (\code{long long},
reduced with the greatest common divisor), then puts all weights over their least common multiple, so a table
holds integers only. The largest intermediate it reaches is $12!$ (the end node's product of node differences for
\code{Central<10>} and the fourth derivative; $12! \cdot 4!$ is a bound that is not reached), far inside
\code{long long}. Every numerator and denominator must fit in 24 bits, which keeps the weights exact in float as
well as in double (the largest numerator is 6,222,216); the generator checks this before narrowing to \code{int},
and a type too narrow to hold a rule's integers, such as half precision under \code{Central<10>}, is rejected by a
\code{static_assert} rather than used with rounded weights. The tests check every table against the moment
conditions with a moment function of their own, and against the published rows of Fornberg\cite{fornberg1988} and
Chapra;\cite[ch.~23]{chapra2015} they also use the exactness on polynomials, with a point and a step that are short
binary fractions so that no rounding occurs, and check that \code{Central<4>} on $x^4$ returns exactly 4 at compile
time and at run time.
\end{sidebar}

\section{The order ladder}\label{sec:stencils-ladder}

\begin{figure}
\sidecaption{Relative error against $h$ for \code{Central<2>} to \code{Central<10>}, $\mathrm{e}^x$ at 1: steeper
truncation slopes, rounding branches that nearly coincide, and minima that slide up the common rounding line as the
best step grows. A higher order buys accuracy by letting the step grow; the minimum falls from about
$\epsmach^{2/3}$ to about $\epsmach^{10/11}$. Data \code{data/step-error-exp.csv}.\label{fig:order-ladder}}
\centering
% F6 fig:order-ladder (Chapter 4, text width). Relative error against the step for Central<2> to Central<10>,
% exp(x) at x = 1 (DS2, data/step-error-exp.csv). The truncation branches steepen with P; the rounding branches
% nearly coincide (S from 1 to 2.3), and the minima slide up the common rounding line eps / h (muted, dashed),
% because a higher order lets the step grow. The five curves have five colours, ordered dark to light with P so
% that they also read in greyscale, and each label takes its curve's colour (the author's release decision: the
% grey ladder was hard to read). The float and caption stay in the chapter.
\definecolor{ladderTwo}{HTML}{2B2B2B}
\definecolor{ladderFour}{HTML}{0B5FAE}
\definecolor{ladderSix}{HTML}{1B8F8A}
\definecolor{ladderEight}{HTML}{6AA84F}
\definecolor{ladderTen}{HTML}{C79A1E}
\begin{tikzpicture}
\begin{axis}[eqtufte, width=95mm, height=45mm,
  xmode=log, ymode=log, log basis x=10, log basis y=10,
  xmin=1e-13, xmax=1, ymin=1e-16, ymax=1,
  xtick={1e-12,1e-10,1e-8,1e-6,1e-4,1e-2,1}, ytick={1e-16,1e-12,1e-8,1e-4,1},
  xlabel={step $h$}, ylabel={relative error}, unbounded coords=jump]
  % the common rounding line eps / h (relative, since f = f' for exp): a guide, not data
  \addplot[black!45, dashed, thin, domain=1e-13:1e-3, samples=2] {2^(-52) / x};
  \node[black!55, font=\scriptsize, anchor=south west, inner sep=1pt] at (axis cs:3e-9, 1.5e-7) {$\epsmach / h$};
  \addplot[ladderTwo] table[col sep=comma, comment chars={\#}, x=h, y=central2] {data/step-error-exp.csv};
  \addplot[ladderFour] table[col sep=comma, comment chars={\#}, x=h, y=central4] {data/step-error-exp.csv};
  \addplot[ladderSix] table[col sep=comma, comment chars={\#}, x=h, y=central6] {data/step-error-exp.csv};
  \addplot[ladderEight] table[col sep=comma, comment chars={\#}, x=h, y=central8] {data/step-error-exp.csv};
  \addplot[ladderTen] table[col sep=comma, comment chars={\#}, x=h, y=central10] {data/step-error-exp.csv};
  % labels on the truncation branches, left of each curve in its colour (the curve is lower to the left)
  \node[ladderTwo, font=\scriptsize, anchor=east, inner sep=1pt] at (axis cs:0.17, 5e-3) {\code{Central<2>}};
  \node[ladderFour, font=\scriptsize, anchor=east, inner sep=1pt] at (axis cs:0.17, 3e-5) {\code{Central<4>}};
  \node[ladderSix, font=\scriptsize, anchor=east, inner sep=1pt] at (axis cs:0.17, 2e-7) {\code{Central<6>}};
  \node[ladderEight, font=\scriptsize, anchor=east, inner sep=1pt] at (axis cs:0.17, 1.6e-9) {\code{Central<8>}};
  \node[ladderTen, font=\scriptsize, anchor=east, inner sep=1pt] at (axis cs:0.17, 1.2e-11) {\code{Central<10>}};
\end{axis}
\end{tikzpicture}

\end{figure}

\Cref{fig:order-ladder} puts the five central rules on one plot. From here to the end of the chapter the function is
$\mathrm{e}^x$ at $x = 1$ rather than the thread, because every derivative of $\mathrm{e}^x$ equals the function,
so the errors of different orders are comparable; the thread returns in \cref{sec:stencils-price}. Three things
to see. On the right, where truncation dominates, each curve is a line whose slope is its order $\aorder$: the
error of \code{Central<4>} falls sixteen times per halving of the step, that of \code{Central<10>} a thousand
times. On the left, where rounding dominates, the curves nearly coincide: the rounding error is $\ampl \epsmach
\abs{f} / h$ with the amplification $\ampl = \sum_k \abs{w_k}$, which grows only from \num{\valAmplCentralTwo} to
\num{\valAmplCentralTen} across the five rules (\cref{tab:weights-central-first}). And the minimum of each V
slides up that common rounding line to the right: a higher order does not lower the rounding branch, it lets the
step grow, and a larger step means a smaller rounding error. The best error falls from about $\epsmach^{2/3}$ for
\code{Central<2>} to about $\epsmach^{10/11}$ for \code{Central<10>}, by the model \num{\valModelCentralTwoFirstDouble}
and \num{\valModelCentralTenFirstDouble}; measured at the default steps the medians over eleven points are
\num{\valExpCentralTwoRelError} and \num{\valExpCentralTenRelError} (\cref{fig:default-step-accuracy}).

\begin{marginfigure}
\centering
% F8 fig:default-step-accuracy (Chapter 4, margin). Relative error at the default step against the accuracy order P
% for Central<2..10> and Forward<1..6>, exp(x), N = 1 (DS7, data/default-step-accuracy.csv): the median over eleven
% points in [0.5, 1], the whiskers the smallest and largest value. Diminishing returns beyond P = 4 or 6 in double.
% The rows with nan (a P the family lacks) are discarded, so the lines join the existing orders.
\begin{tikzpicture}
\begin{axis}[eqtufte, width=38mm, height=34mm,
  ymode=log, log basis y=10, xmin=0.5, xmax=10.5, ymin=1e-16, ymax=1e-7,
  xtick={1,2,...,10}, ytick={1e-16,1e-14,1e-12,1e-10,1e-8},
  xlabel={accuracy order $\aorder$}, ylabel={relative error}, unbounded coords=discard,
  error bars/y dir=both, error bars/y explicit, error bars/error bar style={thin}, error bars/error mark options={rotate=90, mark size=1pt}]
  \addplot[accent, mark=*, mark size=1pt, error bars/error bar style={accent!60}]
    table[col sep=comma, comment chars={\#}, x=P, y=central,
      y error plus expr=\thisrow{central_max}-\thisrow{central}, y error minus expr=\thisrow{central}-\thisrow{central_min}]
    {data/default-step-accuracy.csv};
  \addplot[muted, mark=square*, mark size=1pt, error bars/error bar style={black!35}]
    table[col sep=comma, comment chars={\#}, x=P, y=forward,
      y error plus expr=\thisrow{forward_max}-\thisrow{forward}, y error minus expr=\thisrow{forward}-\thisrow{forward_min}]
    {data/default-step-accuracy.csv};
  \node[accent, font=\tiny, anchor=north, inner sep=1pt] at (axis cs:8, 1e-15) {\code{Central<P>}};
  \node[muted, font=\tiny, anchor=south west, inner sep=1pt] at (axis cs:3.3, 3e-12) {\code{Forward<P>}};
\end{axis}
\end{tikzpicture}

\caption{Relative error at the default step against $\aorder$ for \code{Central<P>} and \code{Forward<P>},
$\mathrm{e}^x$: the median over eleven points in $[0.5, 1]$, the whiskers the smallest and largest value.
Diminishing returns beyond 4 or 6 for double. Data
\code{data/default-step-accuracy.csv}.\label{fig:default-step-accuracy}}
\end{marginfigure}

The returns diminish. \Cref{fig:default-step-accuracy} shows the error at the default step against the order:
going from \code{Central<2>} to \code{Central<4>} gains two decades, from \code{Central<4>} to \code{Central<10>}
less than two more, and the whiskers show that at a single point the error lies up to about
\num{\valDefaultStepSpreadAbove} times above the median and up to \num{\valDefaultStepSpreadBelow} times below it,
by luck in the rounding. In double, \code{Central<4>} is close to as good as it gets for $f'$; the higher orders
earn their keep with wider types, where $\epsmach$ is smaller and $\epsmach^{\aorder/(\aorder+1)}$ has further to
fall (\cref{sec:using-types}).

\section{The price}\label{sec:stencils-price}

What the ladder costs is read off the error model, now written for any rule of order $\aorder$ for $f'$:
\begin{equation}
  E(h) = \errc\, h^{\aorder} \abs{f^{(\aorder + 1)}} + \ampl\, \epsmach \frac{\abs{f}}{h},
  \label{eq:stencils-model}
\end{equation}
with the error constant $\errc$ and the amplification $\ampl$ defined from the weights as above; the concrete
model of \cref{ch:rounding} is the case $\errc = \tfrac16$, $\ampl = 1$. The best step is proportional to
$\epsmach^{1/(\aorder + 1)}$ and the best error to $\epsmach^{\aorder/(\aorder + 1)}$, as in \cref{sb:step-constant}
and \cref{sec:derivations-model}; the default step is $\epsmach^{1/(\aorder + 1)} \scale$ with $\scale = \max(\abs{x}, 1)$,
without the constant that $\errc$ and $\ampl$ would add. On the thread this is visible in the steps themselves:
\code{Central<2>} differentiates $\enth$ with a step of \SI{\valCpDefaultStep}{\kelvin}, \code{Central<4>} with
\SI{\valCpCentralFourStep}{\kelvin} and \code{Central<6>} with \SI{\valCpCentralSixStep}{\kelvin}, and the relative
errors fall from \num{\valCpDefaultRelError} to \num{\valCpCentralFourRelError} and \num{\valCpCentralSixRelError}:
\snippet{04-stencils-orders}{higher-order}

Three prices come with the larger stencil. The first is rounding of the points. \Cref{eq:stencils-model} accounts
for the rounding of the samples $f_k$, but the points $x + n_k h$ are rounded too, each by up to $\epsmach \abs{x}
/ 2$, which moves its sample by about that times $f'$. Over the stencil this adds a third term,
$(\ampl / 2)\, \epsmach \abs{x f'} / h$, of the same order as the second. The two-node rules escape it: they divide
by the distance between the points actually evaluated rather than by $2h$, which removes the point rounding exactly
(\cref{sb:actual-distance}). A rule with more nodes cannot. Dividing by the distance between its outer nodes would
not remove their rounding but quadruple it, since their weight in \code{Central<4>} is $1/12$ and the distance
carries the rounding with weight $4/12$, and it would leave the inner nodes' rounding untouched; so the rule divides
by a fixed multiple of $h$ and keeps the third term, which the tests allow for (\cref{sec:guarantees-verification}).

The second price is smoothness. The error term in \cref{eq:stencils-central-four} holds only if $f$ has
$\aorder + \dorder$ continuous derivatives across the whole stencil, which spans $K h$ on each side of $x$ with $K$
the largest node, and both $K$ and the default step grow with the order: on the thread \code{Central<6>} samples
up to \SI{\valCpCentralSixStep}{\kelvin} times three on either side of $T_0$. A kink or a range break anywhere in
that span gives a wrong value without any failure, and a higher order makes it worse, not better; \cref{sec:using-pitfalls}
measures this on the thread's own range break. The third price is evaluations: $\aorder + \dorder - 1$ per call
(\cref{tab:stencils-nodes}), which inside an iteration is the whole cost of the derivative.

\begin{sidebar}[label=sb:richardson-disguise]{\code{Central<4>} is Richardson extrapolation in disguise}
Let $\Dc{h} = (f(x + h) - f(x - h)) / 2h$ be the central difference. Its error expansion has even powers only,
$\Dc{h} = f'(x) + a_2 h^2 + a_4 h^4 + \cdots$ (\cref{sb:remainder}), so halving the step divides the leading term
by 4, and one combination removes it:
\[
  \frac{4 \Dc{h/2} - \Dc{h}}{3} = f'(x) - \frac{a_4}{4} h^4 + \cdots .
\]
Writing out the four samples with $h' = h/2$ gives
$(8 f(x + h') - 8 f(x - h') - f(x + 2h') + f(x - 2h')) / (12 h')$, which is \cref{eq:stencils-central-four}
exactly. Richardson extrapolation at one fixed level is a stencil: the same samples, the same weights, the same
error, and no information about how large that error is. What Ridders' method (\cref{ch:ridders}) adds is not a
better formula but a decision: it keeps halving and extrapolating until the entries stop changing by more than
rounding, and reads an error estimate off the changes. The price is that the number of samples is decided at run
time.
\end{sidebar}

\section{One-sided rules}\label{sec:stencils-one-sided}

\begin{figure}
\sidecaption{The nodes of \code{Central<4>}, \code{Forward<2>} and \code{Backward<2>} for $\dorder = 1$ and
$\dorder = 2$ on a line: node 0 filled where the rule uses it. A central rule for an odd derivative does not sample
$x$ itself. A diagram without data.\label{fig:stencil-nodes}}
\centering
% F9 fig:stencil-nodes (Chapter 4, text width). The nodes of Central<4>, Forward<2> and Backward<2> for N = 1 and
% N = 2 on a line (a diagram, no data): a filled dot is a node the rule samples, in units of h from x. Node 0 (the
% point x itself) is in accent, filled where the rule uses it and hollow where it does not. From deriv.hpp:
% Central<P> uses the P + N - 1 integers around 0, without 0 for odd N and with 0 for even N; Forward<P> uses
% 0 ... P + N - 1; Backward<P> is Forward<P> mirrored.
\begin{tikzpicture}[x=7.5mm, y=5.5mm, font=\footnotesize,
  used/.style={fill=black!75, draw=none, circle, inner sep=0pt, minimum size=4pt},
  centre used/.style={fill=accent, draw=none, circle, inner sep=0pt, minimum size=4pt},
  centre unused/.style={fill=white, draw=accent, thin, circle, inner sep=0pt, minimum size=4pt}]
  % one row per rule: label, the line with ticks, the sampled nodes; the centre node drawn last
  \foreach \row/\name/\order/\nodes/\usesCentre in {
      0/Central<4>/1/{-2,-1,1,2}/0,
      1/Central<4>/2/{-2,-1,0,1,2}/1,
      2/Forward<2>/1/{0,1,2}/1,
      3/Forward<2>/2/{0,1,2,3}/1,
      4/Backward<2>/1/{-2,-1,0}/1,
      5/Backward<2>/2/{-3,-2,-1,0}/1} {
    \draw[black!55, thin] (-3.6, -\row) -- (3.6, -\row);
    \foreach \k in {-3,...,3} \draw[black!55, thin] (\k, -\row) ++(0, -0.12) -- ++(0, 0.24);
    \node[anchor=east] at (-4.0, -\row) {\texttt{\name}, $\dorder = \order$};
    \foreach \k in \nodes { \ifnum\k=0 \else \node[used] at (\k, -\row) {}; \fi }
    \ifnum\usesCentre=1 \node[centre used] at (0, -\row) {}; \else \node[centre unused] at (0, -\row) {}; \fi
  }
  % node labels under the last row
  \foreach \k in {-3,...,3} \node[anchor=north, font=\scriptsize, black!70] at (\k, -5.3) {$\k$};
  \node[anchor=north, font=\scriptsize, black!70] at (0, -5.95) {node $n_k$: the sample is at $x + n_k h$};
  % the one thing to notice: the centre is not sampled by Central<4> for f'
  \node[accent, font=\scriptsize, anchor=west] at (3.8, 0) {$x$ not sampled};
  \node[accent, font=\scriptsize, anchor=west] at (3.8, -1) {$x$ sampled};
\end{tikzpicture}

\end{figure}

A central rule samples both sides of $x$. Where $f$ exists on one side only, at the first temperature of a
correlation, at a phase boundary, at $t = 0$ of a simulation, the samples on the other side are either
unavailable or meaningless. The one-sided rules put all their nodes on one side: \code{Forward<P>} uses the nodes
$0, 1, \ldots, \aorder$ for $f'$ and never evaluates below $x$; \code{Backward<P>} is its mirror image, the nodes
$-\aorder, \ldots, 0$, with the weights reversed and multiplied by $(-1)^{\dorder}$, and never evaluates above
$x$ (\cref{fig:stencil-nodes}). Both include $x$ itself, so they also serve at the upper end of a range, where
\code{Backward<P>} stays inside it, and for a function that is cheap to evaluate on one side only.

% The margin figure is anchored before the paragraph that cites it so that its caption clears the margin caption of
% fig:stencil-nodes at the foot of the same page (stage-6 layout check).
\begin{marginfigure}
\centering
% F7 fig:one-sided (Chapter 4, margin). Forward<2> against Central<2> and Forward<4> against Central<4>, exp(x) at
% x = 1 (DS2): the same truncation slopes, higher rounding floors for the one-sided rules (larger S). One-sided rules
% in accent, central rules muted.
\begin{tikzpicture}
\begin{axis}[eqtufte, width=35mm, height=34mm,
  xmode=log, ymode=log, log basis x=10, log basis y=10,
  xmin=1e-13, xmax=1, ymin=1e-16, ymax=10,
  xtick={1e-12,1e-8,1e-4,1}, ytick={1e-16,1e-12,1e-8,1e-4,1},
  xlabel={step $h$}, ylabel={relative error}, unbounded coords=jump]
  \addplot[muted] table[col sep=comma, comment chars={\#}, x=h, y=central2] {data/step-error-exp.csv};
  \addplot[muted] table[col sep=comma, comment chars={\#}, x=h, y=central4] {data/step-error-exp.csv};
  \addplot[accent] table[col sep=comma, comment chars={\#}, x=h, y=forward2] {data/step-error-exp.csv};
  \addplot[accent] table[col sep=comma, comment chars={\#}, x=h, y=forward4] {data/step-error-exp.csv};
  % In the margin width the two rules of one order lie on top of each other along the truncation branch, so each
  % pair carries one order label, and the family is given by the colour (the key in the free top-left corner).
  \node[accent, font=\tiny, anchor=north west, inner sep=1pt] at (axis cs:1.5e-13, 5) {\code{Forward<P>}};
  \node[muted, font=\tiny, anchor=north west, inner sep=1pt] at (axis cs:1.5e-13, 0.15) {\code{Central<P>}};
  \node[font=\tiny, anchor=east, inner sep=1pt] at (axis cs:0.004, 1e-5) {$\aorder = 2$};
  \node[font=\tiny, anchor=east, inner sep=1pt] at (axis cs:0.004, 3e-10) {$\aorder = 4$};
\end{axis}
\end{tikzpicture}

\caption{\code{Forward<2>} against \code{Central<2>} and \code{Forward<4>} against \code{Central<4>}: the same
slopes, higher floors (a larger $\ampl$). Data \code{data/step-error-exp.csv}.\label{fig:one-sided}}
\end{marginfigure}

Each point buys one order. \code{Forward<P>} has $\aorder + 1$ weights and no symmetry to help, so the
$\aorder + 1$ moment conditions $\mom{0}, \ldots, \mom{\aorder}$ fix them and $\mom{\aorder + 1}$ is the first
that does not vanish: order $\aorder$ from $\aorder + 1$ points. A central rule gets its even moments for nothing
from antisymmetry, which is why it reaches the same order with one point fewer, and why its orders are even.
The one-sided rules pay in their constants. The amplification $\ampl$ of \code{Forward<1>} to \code{Forward<6>}
runs from \num{\valAmplForwardOne} to \num{\valAmplForwardSix} against \num{\valAmplCentralTwo} to
\num{\valAmplCentralTen} for the central rules, and the error constant of \code{Forward<2>} is
\num{\valErrcForwardTwo} against \num{\valErrcCentralTwo} (\cref{tab:weights-forward-first}). The weights
alternate in sign and grow, so the samples cancel more and the rounding is amplified more.
\Cref{fig:one-sided} shows the effect: the same truncation slopes as the central rules of the same order, with
the rounding branch shifted up. At the default step on $\mathrm{e}^x$, \code{Forward<2>} reaches
\num{\valExpForwardTwoRelError} where \code{Central<2>} reaches \num{\valExpCentralTwoRelError}, and
\code{Forward<6>} \num{\valExpForwardSixRelError} where \code{Central<6>} reaches \num{\valExpCentralSixRelError}
(\cref{fig:default-step-accuracy}). The one-sided rules stop at order 6, a design decision discussed with the
central cap in \cref{sec:stencils-choosing}.

The thread shows why the edge matters. Its Shomate fit is valid from \SI{500}{\kelvin} to \SI{2000}{\kelvin}, but
the formula, a polynomial plus a $1/T$ term, evaluates without complaint at \SI{499}{\kelvin} and returns a
plausible number that describes nothing: evaluable is not valid. A central rule at \SI{500}{\kelvin} samples
below the range and returns a derivative of the extrapolated formula, with no sign that anything is wrong. The
remedy is to make the domain explicit. The wrapper below returns NaN outside the fit's range, so that any rule
sampling there fails with \code{non_finite} instead of succeeding; the one-sided rules never sample there and
succeed, \code{Forward<2>} with its three points and \code{Forward<4>} with five:
\needspace{14\baselineskip}%
\snippet{04-stencils-edge}{validity-edge}
Measured, \code{Forward<2>} is \num{\valEdgeForwardTwoRelError} off the closed form and \code{Forward<4>}
\num{\valEdgeForwardFourRelError}. Ridders' method, which restarts below a step at which $f$ is not finite
(\cref{ch:ridders}), does not help exactly at an edge: with the same wrapper at \SI{500}{\kelvin} every level
samples below the range, each restart halves the step in vain, and the result is \valEdgeRiddersStatus{} after
\num{\valEdgeRiddersEvals} evaluations, the level cap; near such an edge rather than at it, as for $\log x$ close
to 0, it recovers (\cref{sec:ridders-trust}). Below \SI{500}{\kelvin} the published enthalpy of nitrogen is a
different fit, and the two do not meet, which makes the junction the range break of \cref{sec:using-pitfalls}.

\section{Choosing P}\label{sec:stencils-choosing}

As a rule: \code{Central<2>} or \code{Central<4>} for double inside an iteration, where two or four evaluations buy
a relative error of about \num{\valExpCentralTwoRelError} or \num{\valExpCentralFourRelError} with the default
step; \code{Central<6>} and above only for wide types, where $\epsmach$ is small enough for the order to pay
(\cref{sec:using-types} uses \code{Central<6>} with fifty digits); \code{Forward<2>} or \code{Forward<4>}, or
their mirrors, exactly at an edge and nowhere else, since their constants are larger. For a value to report with
an error estimate, use Ridders. The caps are design decisions recorded in the design record of the repository
(\code{docs/PHASE2_API.md}), not measured limits. The central rules stop at order 10: by the model's exponent
alone, a rule of order 12 would gain about \num{\valCapGainCentralTwelveFirstDouble} digit for $f'$ in double and
\num{\valCapGainCentralTwelveFirstMultiprecision} in a 50-digit type (\num{\valCapGainCentralTwelveFourthDouble} and
\num{\valCapGainCentralTwelveFourthMultiprecision} for the fourth derivative), and the library stops there. The
one-sided rules stop at order 6, where the amplification, already \num{\valAmplForwardSix} for $f'$, grows about
as fast as the exponent pays (\cref{sec:guarantees-limits}). The decision table by situation is
\cref{tab:method-choice}, and every rule's nodes, weights and constants are in \cref{ch:tables}.
